\section{Materiales y sombreadores (Shaders)}
Para la visualización de la geometría se optó por un modelo de renderizado no basado en física (\textit{Non-PBR}, 
por sus siglas en inglés), en virtud de que el objetivo central es facilitar la interpretación y el análisis cuantitativo d
e los tensores de la red. Al prescindir del renderizado fotorrealista, el sistema prioriza la fidelidad cromática y la 
legibilidad de los datos sobre efectos estéticos complejos, adoptando el estándar habitual de la visualización científica. 
Con este propósito, todos los sombreadores del sistema se fundamentan en nodos de emisión pura (\textit{Emission Shaders}).
\medskip

En computación gráfica, los materiales emisivos representan superficies que proyectan luz propia. Como resultado, estos 
sombreadores omiten el cómputo de sombras, oclusiones o reflexiones especulares dependientes de fuentes de luz externas, 
garantizando que los valores de color se desplieguen de forma homogénea e invariable respecto a las fuentes de iluminación
de la escena. Dicha propiedad convierte a los materiales emisivos en la solución idónea para presentar texturas y 
mapas escalares bajo su valor cromático directo (color real o \textit{albedo}).
\medskip

El material \textbf{Volume Slice} se encarga del sombreado de las rebanadas volumétricas. Este material recibe como atributo 
geométrico un valor escalar normalizado en el intervalo $[0, 1]$, y mediante un nodo de rampa de color (\textit{Color Ramp}) 
mapea dicho escalar a un gradiente predefinido. El valor transferido corresponde a la magnitud de activación del tensor en 
la celda evaluada, generando así un mapa de calor (\textit{heatmap}) que refleja la intensidad de las activaciones intermedias.
\medskip

El material \textbf{Simplified Volume} se emplea para el sombreado de las capas en su modo de representación compacta. A 
partir de un color base asignado por el nodo de geometría (utilizado para clasificar visualmente el tipo de capa u operación), 
el sombreador aplica un gradiente vertical sutil para dotar de profundidad y contraste al cuerpo geométrico simplificado.
\medskip

Por su parte, el material \textbf{Volume BB} se destina a los delimitadores de volumen (\textit{bounding boxes}), las aristas 
de conexión entre capas y las geometrías de los kernels. Este material aplica una emisión blanca uniforme para establecer un 
contraste visual nítido frente al fondo neutro del espacio de trabajo y las superficies texturizadas de las rebanadas.
\medskip

El material \textbf{Box Class} permite mapear proceduralmente las etiquetas de clase a un espectro de color identificativo. 
De manera análoga a \textbf{Volume Slice}, recibe como parámetro un valor normalizado en el rango $[0, 1]$, correspondiente 
al identificador numérico de la clase (\textit{class\_id}) escalado por el número total de categorías. Con este valor se 
modula directamente el matiz (\textit{hue}) en el espacio de color HSV del material, garantizando una codificación cromática 
determinista y diferenciable por cada clase detectada.
\medskip

Finalmente, el material \textbf{Input Image} proyecta la imagen de entrada de la red sobre el plano focal mediante un mapeo 
de textura con emisión directa, sirviendo como marco de referencia visual para la superposición de las cajas delimitadoras.